\begin{exercisechunk}
\item \textbf{The variation norm of a univariate ReLU network.}
\label[exercise]{ex:l15-relu-spline-variation}
\exerciseprofile{3}{3}{3}{\exproof\quad\excalculation\quad\exinterpretation}{%
Derive the function-space regularizer induced by squared parameter-norm
control and relate its atomic measures to the knots of a linear spline.}
Consider finite two-layer ReLU networks
\[
 h_\theta(x)=c+\sum_{j=1}^m a_j[w_jx+b_j]_+,
 \qquad
 C(\theta)=\frac12\sum_{j=1}^m(a_j^2+w_j^2),
\]
where the biases are unregularized. For a continuous function \(f:\R\to\R\),
define its unbounded-width representation cost by
\begin{align*}
 \cR(f)=\inf\biggl\{
 &\|\mu_+\|_{\TV}+\|\mu_-\|_{\TV}:\\
 &f(x)=c+\int_\R[x-b]_+\,d\mu_+(b)
          +\int_\R[b-x]_+\,d\mu_-(b)
 \text{ for some }c\in\R\biggr\}.
\end{align*}
Assume below that \(D^2f\) is a compactly supported finite signed measure and
that the limiting slopes \(f'(-\infty)\) and \(f'(+\infty)\) exist. You may use
\(D^2[x-b]_+=D^2[b-x]_+=\delta_b\) in the sense of distributions.
\begin{enumerate}[(a),beginpenalty=10000]
\item Fix one nonconstant unit \(a[wx+b]_+\). Show that its function is
unchanged when \((a,w,b)\) is replaced by
\((a/r,rw,rb)\) for \(r>0\). Prove
\[
 \inf_{r>0}\frac12\left(\frac{a^2}{r^2}+r^2w^2\right)
 =|aw|.
\]
Explain why the sign of \(w\) then produces one of the two normalized atoms
\([x-t]_+\) and \([t-x]_+\).

\item For any measure representation in the definition of \(\cR(f)\), prove
\[
 D^2f=\mu_++\mu_-,
 \qquad
 f'(-\infty)+f'(+\infty)
 =\mu_+(\R)-\mu_-(\R).
\]
Set \(B=\|D^2f\|_{\TV}\) and
\(t=f'(-\infty)+f'(+\infty)\). Use the two identities and the triangle
inequality to prove
\[
 \cR(f)\geq\max\{B,|t|\}.
\]

\item Prove the matching upper bound and conclude that
\[
 \cR(f)=\max\left\{
 \|D^2f\|_{\TV},
 \left|f'(-\infty)+f'(+\infty)\right|
 \right\}.
\]
Your construction must also cover the case \(B=0\).

\item Let \(s:\R\to\R\) be a linear spline with ordered knots
\(x_1<\cdots<x_N\). Let its slopes on the successive intervals be
\(\ell_0,\ell_1,\ldots,\ell_N\). Prove
\[
 D^2s=\sum_{i=1}^N(\ell_i-\ell_{i-1})\delta_{x_i}
\]
and hence
\[
 \cR(s)=\max\left\{
   \sum_{i=1}^N|\ell_i-\ell_{i-1}|,
   |\ell_0+\ell_N|
 \right\}.
\]
If the knots are ordered sample locations and \(s(x_i)=y_i\), identify the
interior slopes \(\ell_1,\ldots,\ell_{N-1}\) and write the remaining
minimization over the two exterior slopes explicitly.

\item The result cited in the bibliographical notes proves that a
minimum-cost interpolant over the full function class can be chosen from the
linear splines in part (d). Explain what parts (a)--(d) establish about the
relation between neural parameter regularization and splines. Explain why
the existence of a linear-spline minimizer does not imply that every
minimum-cost interpolant is that spline.
\end{enumerate}
\begin{hints}
For part (c), first suppose \(B>0\). Write
\(D^2f=q\,|D^2f|\), where \(|q|=1\) almost everywhere, and introduce a
signed measure \(\eta=(t/B)|D^2f|\). Seek \(\mu_+\) and \(\mu_-\) whose sum
is \(D^2f\) and whose difference is \(\eta\). Treat \(B=0\) separately.
\end{hints}
\end{exercisechunk}
